\section{Operational Game Semantics}
\label{s:AOGS}
 
We introduce the representation of $\LasV$ in OGS. 
The LTS produced by the embedding of a term intends to capture the possible interactions 
between this term and its environment.
Values exchanged between the term and the environment are represented by free variables, 
ranged over by $x,y,z$. 
Continuations (i.e., evaluation contexts) are also
represented by names, called \emph{continuation names} and ranged over by  $p,q,r$.

\emph{Actions} $\act$ have been introduced in Section~\ref{s:nota}. 
In OGS, we have five kinds of (visible) actions:
\begin{itemize}
  \item \emph{Player Answers} (PA), $\ansP{p}{x}$,
  and \emph{Opponent Answers} (OA),  $\ansO{p}{x}$,
  that exchange a variable $x$ through a continuation name $p$;
  \item \emph{Player Questions} (PQ), $\questP{x}{y,p}$,
  and \emph{Opponent Questions} (OQ), $\questO{x}{y,p}$, 
  that exchange a variable $y$ and a continuation name $p$ through a variable $x$;
  \item \emph{Initial Opponent Questions} (IOQ), $\questOinit{p}$, that
  introduce the initial continuation name $p$.
\end{itemize}
 
\begin{remark}
 \label{rmk:justified-sequence}
The denotation of terms is usually represented in game semantics using 
the notion of \emph{pointer structure} rather than traces.
A pointer structure is defined as a sequence of \emph{moves}, 
together with a pointer from each move (but the initial one) to a previous move 
that  `justifies' it.
Taking a trace $\tr$, one can reconstruct this pointer structure in the following way:
an action $\act$ is \emph{justified} by an action $\act'$
if the free name of $\act$ is bound by $\act'$ in $\tr$ (here we are taking advantage of
the `freshness' convention on the bound names of traces, Remark~\ref{r:bn}).
\end{remark}


\emph{Environments}, ranged over by $\gamma$, 
maintain the association from names to values and evaluation contexts, and  
 are partial maps with finite domain. A single mapping is either of the form 
$\pmap{x}{V}$
(the variable $x$ is mapped onto the  value $V$), or 
$\pmap{p}{(E,q)}$ 
(the continuation name $p$ is mapped onto the 
pair of the  evaluation context $E$ and the continuation $q$).
We write $\emptymap$ for the empty environment,
and $\gamma \cdot \gamma'$ for the union of two such partial maps,
provided that their domains are disjoint.


There are two main kinds of \emph{configurations} $\FF$:
\emph{active configurations} $\conf{M,p,\gamma,\phi}$ and 
\emph{passive configurations}
$\conf{\gamma,\phi}$, where $M$ is a term, $p$ a continuation name, 
$\gamma$ an environment and $\phi$ a set of names called its \emph{name-support},
corresponding to the set of all names that has been exchanged so far.
 

 \begin{figure*}
 \[
  \begin{array}{l|llll}
   (P\tau) & \conf{M,p,\gamma,\phi} & \arrA{\ \tau \ } & 
     \conf{N,p,\gamma,\phi} & \text{ when } M \redv N\\
%
   (PA) & \conf{V,p,\gamma,\phi} & \arrA{\ansP{p}{x}} & 
     \conf{\gamma \cdot \pmap{x}{V},\phi \uplus \{x\}} \\
%
   (PQ) & \conf{E[xV],p,\gamma,\phi} & \arrA{\questPV{x}{y}{q}} & 
     \multicolumn{2}{l}{\conf{\gamma \cdot \pmap{y}{V}\cdot \pmap{q}{(E,p)},\phi 
       \uplus \{y,q\}}} \\
%
   (OA) & \conf{\gamma\cdot \pmap{q}{(E,p)},\phi} & \arrA{\ansO{q}{x}} & 
     \conf{E[x],p,\gamma,\phi \uplus \{x\}}\\
%
   (OQ) & \conf{\gamma,\phi} & \arrA{\questOV{x}{y}{p}} & 
     \conf{V y,p,\gamma,\phi \uplus \{y,p\}} & \text{ when } \gamma(x) = V\\
%
   (IOQ) & \initconf{\phi}{M} & \arrA{\questOinit{p}} &
     \conf{M,p,\emptymap,\phi \uplus \{p\}}
  \end{array}\]
 \caption{The  LTS for the  Alternating OGS (\AOGS)}
 \label{fig:seq-ogs-lts}
\end{figure*}

The LTS is introduced in Figure~\ref{fig:seq-ogs-lts}.
It is called \emph{Alternating}, since, forgetting the P$\tau$ transition,
it is bipartite between active configurations that perform Player actions and $\tau$-transitions,
and passive configurations that perform Opponent actions.  
Accordingly, we call Alternating the resulting OGS, abbreviated \AOGS. 
In the OA rule, $E$ is ``garbage-collected''  from $\gamma$,
a behavior corresponding to \emph{linear continuations}.

 
The term of an active configuration determines the next transition performed.
First, the term needs to be reduced, using the rule (P$\tau$).
When the term is a value $V$, a Player Answer (PA) is performed,
providing a fresh variable $x$ to Opponent, while $V$ is stored in $\gamma$ at position $x$. 
Freshness is enforced using the disjoint union $\uplus$.
When the term is a callback $E[xV]$, with $p$ the current continuation name,
a Player Question (PQ) at $x$ is performed, providing two fresh names $y,q$
to Opponent, while storing $V$ at $y$ and $(E,p)$ at $q$ in $\gamma$.

 
On passive configurations, Opponent has the choice to perform different actions.
It can  perform an Opponent Answer (OA) by interrogating an evaluation context 
$E$ stored in $\gamma$. For this, Opponent provides a fresh variable $x$ that is plugged into the hole of $E$,
while the continuation name $q$ associated with $E$ in $\gamma$ is restored.
Opponent may also  perform  an Opponent Question (OQ),
by interrogating a value $V$ stored in $\gamma$. 
For this, Opponent provides a fresh variable $y$ as an argument to $V$.


To build the denotation of a term $M$, we introduce an \emph{initial
configuration} associated with it, 
written $\initconf{\phi}{M}$, with $\phi$ the set of free variables we start with.
When this set is taken to be the free variables of $M$, we simply write it as $\conf{M}$.
In the initial configuration, the choice of the continuation
name $p$ is made by performing an Initial Opponent question
(IOQ). So initial configurations are passive.

From an environment $\gamma$, we extract its \emph{continuation structure}
which is defined
as the relation over continuation names 
\[\cs(\gamma) \defeq \{(p,q) \mid p \in \dom{\gamma} \text{ and } \gamma(p) = (\_,q)\}\]
As we will see, the continuation structure indicates
on which continuation name $q$ Player should answer 
following an Opponent answer over $p$.
Continuation names of $\phi$ that are either
the toplevel continuation name -- for active configurations --
or that are in the continuation structure of $\gamma$ are called \emph{available}.
So continuation names that are not available are the ones that
have already been used to perform an answer.
This distinction is needed since we consider linear continuations here.
Variables are always considered to be available since there is
no constraint on the number of times they can be interrogated.

Names in $\dom{\gamma}$ are called \emph{P-names}, 
and those in $\phi \backslash \dom{\gamma}$ that are available are called \emph{O-names}.
So we obtain a \emph{polarity function} $\Label_{\FF}$ 
associated to $\FF$, 
defined as the partial maps
from $\phi$ to $\{O,P\}$  mapping names to their polarity.
 %
 In the following, we  only consider \emph{valid configurations},
 for which:
 \begin{itemize}
  \item $\dom{\gamma} \subseteq \phi$
  \item $\FV(M),p$ are O-names;
  \item for all $a \in \dom{\gamma}$, the names appearing in $\gamma(a)$ are O-names.
 \end{itemize}